% !TeX program = xelatex
% Сборка (из папки submission): xelatex report.tex — два раза (ссылки и номера).
% Отчёт склеен из sections/*.tex: каждый правит только свои файлы.
% Красный [TODO: ...] — незаполненное место; перед сдачей их быть не должно.
\documentclass[12pt,a4paper]{article}
\usepackage{fontspec}
\usepackage{polyglossia}
\setdefaultlanguage{russian}
% Шрифты — как в шаблоне преподавателя (report_template.tex).
% Если Times New Roman / Courier New не установлены (типично для Linux),
% берётся метрически совместимая замена Liberation Serif / Liberation Mono
% (пакет fonts-liberation), чтобы сборка не падала.
% Для точного совпадения с шаблоном: ttf-mscorefonts-installer (см. README.md).
\IfFontExistsTF{Times New Roman}{%
  \setmainfont{Times New Roman}%
  \setmonofont{Courier New}%
  \newfontfamily\cyrillicfonttt[Script=Default]{Courier New}%
}{%
  \setmainfont{Liberation Serif}%
  \setmonofont{Liberation Mono}%
  \newfontfamily\cyrillicfonttt[Script=Default]{Liberation Mono}%
}
\usepackage{amsmath,amssymb,graphicx,array}
\usepackage[margin=2cm]{geometry}
\usepackage{xcolor,listings}
\usepackage[unicode,hidelinks]{hyperref}
\setlength{\parindent}{0pt}
\setlength{\parskip}{6pt}
\setlength{\emergencystretch}{3em}
\renewcommand{\arraystretch}{1.25}
\newcommand{\taskhead}[1]{\par\textbf{#1}\par}
% Сохранять ASCII-дефис U+002D при копировании кода из PDF.
\lstset{literate={-}{{\char45}}1}
\lstdefinestyle{commands}{
  basicstyle=\ttfamily\small,
  backgroundcolor=\color{gray!10},
  frame=single,framerule=0pt,framesep=6pt,
  xleftmargin=6pt,xrightmargin=6pt,
  breaklines=true,columns=fullflexible,keepspaces=true,
  showstringspaces=false
}
% Красная пометка о незаполненном месте.
\newcommand{\todo}[1]{\textcolor{red}{[TODO: #1]}}
% Рисунок с защитой от отсутствующего файла: пока файла нет — рамка-заглушка,
% сборка не ломается. Пример:
%   \taskfigure[0.75]{figures/2-2-error.pdf}{Подпись}{fig:2-2-error}
\newcommand{\taskfigure}[4][0.75]{%
  \begin{figure}[htbp]\centering
  \IfFileExists{#2}{\includegraphics[width=#1\linewidth]{#2}}{%
    \fbox{\parbox[c][0.2\textheight][c]{0.85\linewidth}{\centering
      \todo{рисунок ещё не добавлен}\\[6pt]\texttt{\detokenize{#2}}}}}%
  \caption{#3}\label{#4}%
  \end{figure}}

\begin{document}
\sloppy
% ======================================================================
% Титульный лист. Ответственный: A — Плешнев Арсений.
% Требования: main.tex, строки 1847–1850 (название, тема, ФИО команды,
% номер группы, номер варианта; цель своими словами; ОС и версии библиотек).
% Образец оформления: report_template.tex, строки 34–76.
% ======================================================================
\begin{titlepage}
{\large\bfseries Отчёт по лабораторной работе\par}
\textbf{По дисциплине:} Математический анализ (3 семестр)

\textbf{На тему:} Численное вычисление кратных интегралов.
Исследование и оптимизация динамической системы
\vspace{1cm}

\textbf{Команда №1}

\textbf{Вариант № 3}

\textbf{Выполнили} студенты гр. № \todo{A: номер группы}:

\textit{Плешнев Арсений Вадимович, Попадюк Евгений Алексеевич,
Дучин Кирилл Дмитриевич, Филиппов Андрей Алексеевич}

\textbf{Дата выполнения:} \todo{A: дата}
\vspace{0.7cm}

\textbf{Цель работы}

\todo{A: сформулировать цель своими словами (2–3 предложения, не копировать условие).}
\vspace{0.5cm}

\textbf{Условия выполнения}

\todo{A: ОС, версии Python, NumPy, Matplotlib; сборка XeLaTeX.
Команда для вывода версий — в README.md. Заполнить по машине,
на которой получены итоговые результаты.}
\end{titlepage}

% ======================================================================
% Порядок выполнения. Ответственный: A (собирает строки от всех).
% Требования: main.tex, строки 1851–1853 — краткая последовательность
% фактически выполненных действий разделов 1–3; какие пункты со звёздочкой выполнены.
% Образец: report_template.tex, строки 78–88.
% Каждый владелец присылает A одну строку про свою часть (или правит свою строку через PR).
% ======================================================================
\taskhead{Порядок выполнения}
\begin{enumerate}
\item[0] \todo{A: подготовка окружения (Python, библиотеки, XeLaTeX).}
\item[1.1] \todo{D: что сделано — аналитика, обе квадратуры, таблицы, графики.}
\item[1.2] \todo{D: что сделано — аналитика, замена переменных, якобиан, численный расчёт.}
\item[2.1] \todo{D: что сделано.}
\item[2.2] \todo{D: что сделано.}
\item[2.3] \todo{D: что сделано.}
\item[3.1] \todo{A: что сделано.}
\item[3.2] \todo{B: что сделано.}
\item[3.3] \todo{B: что сделано.}
\item[3.4] \todo{A: что сделано.}
\item[3.5] \todo{A/B/C: численные эксперименты 3.5.1–3.5.5.}
\item[3.6] \todo{C: что сделано.}
\end{enumerate}
\todo{A: указать, какие пункты со звёздочкой (3.1.6*, 3.2.5*) выполнены, или что они не выполнялись.}


\section{Численное вычисление двойных интегралов}
% ======================================================================
% Раздел 1. Двойные интегралы (задачи 1.1 и 1.2).
% Ответственный: D — Филиппов Андрей.
%
% Что сдать (main.tex: порядок 249–283, разобранный пример 288–387;
% требования к отчёту 1854–1861; правила точности 1948–1999;
% пример таблицы и графика ошибки 2001–2083):
%   1.1  Аналитически вычислить I1 по прямоугольнику [0,a]x[0,b];
%        реализовать метод средних прямоугольников и метод трапеций;
%        вычислить I1 при n = 4, 8, 16, 32, 64 по каждой координате.
%   1.2  Аналитически вычислить I2 по треугольнику D; перейти к единичному
%        квадрату (x=a*xi, y=b*(1-xi)*eta), найти якобиан и вычислить
%        преобразованный интеграл численно обоими методами.
%   Для обеих задач: изображения областей, таблицы сравнения с аналитикой,
%   графики абсолютной ошибки от шага, вывод об изменении ошибки при удвоении n.
%   Нулевые ошибки на логарифмическую шкалу не наносить; если метод даёт точное
%   значение с учётом округления — указать это в отчёте.
%
% Параметры варианта 3 (приложение А): a=2, b=3.
% ВНИМАНИЕ: a,b,c приложения А (интегралы) и приложения Б (модель) — разные величины.
% Шаги: h_x=a/n, h_y=b/n; на графике h=max(h_x,h_y); после замены h=1/n.
%
% Код квадратур — в solution.py (блок D), вызовы — task_1_1/task_1_2
% в experiments.py (запуск: python experiments.py 1.1).
% Рисунки: figures/1-1-*.pdf, figures/1-2-*.pdf (сохранять из experiments.py).
% ======================================================================

\subsection{Интеграл по прямоугольной области}

\taskhead{Исходные данные и аналитическое решение}
Для варианта 3 заданы $a=2$, $b=3$. Рассматривается интеграл
\[
I_1=\int_0^2\int_0^3(x^2+xy+y^2)\,dy\,dx.
\]
\todo{D: вычислить $I_1$ аналитически — повторное интегрирование с промежуточными шагами (пример: main.tex, строки 292–304).}

\taskhead{Программный результат и сравнение}
\todo{D: листинг вызова своих функций из solution.py и полученные значения (образец: \texttt{report\_template.tex}, строки 138–149).}
Аналитическое значение $I_1$: \todo{D: подставить}. Шаги сетки $h_x=a/n$, $h_y=b/n$;
для графика используется $h=\max(h_x,h_y)$. Ошибки $\Delta_n=|I_n-I_1|$,
$\delta_n=\Delta_n/|I_1|$ вычисляются до округления результатов.
\begin{table}[htbp]
\centering\small
\caption{Сравнение численного интеграла с аналитическим значением $I_1$}
\label{tab:1-1}
\begin{tabular}{|l|r|r|r|r|r|}\hline
Метод & $n$ & $h$ & $I_n$ & $\Delta_n$ & $\delta_n$\\\hline
Ср. прямоугольники & 4 & — & — & — & —\\
Ср. прямоугольники & 8 & — & — & — & —\\
Ср. прямоугольники & 16 & — & — & — & —\\
Ср. прямоугольники & 32 & — & — & — & —\\
Ср. прямоугольники & 64 & — & — & — & —\\
Трапеции & 4 & — & — & — & —\\
Трапеции & 8 & — & — & — & —\\
Трапеции & 16 & — & — & — & —\\
Трапеции & 32 & — & — & — & —\\
Трапеции & 64 & — & — & — & —\\
\hline\end{tabular}
\end{table}

\taskhead{Графическая проверка}
\taskfigure[0.6]{figures/1-1-region.pdf}{Область интегрирования — прямоугольник $[0,a]\times[0,b]$ (равный масштаб осей)}{fig:1-1-region}
\taskfigure[0.75]{figures/1-1-error.pdf}{Абсолютная ошибка обоих методов от шага сетки (логарифмические оси)}{fig:1-1-error}
\todo{D: пояснить графики: как меняется ошибка при удвоении $n$ и почему.}

\taskhead{Вывод}
\todo{D: вывод о сходимости обоих методов; если метод даёт точное значение с учётом округления — указать это.}

\subsection{Интеграл по треугольной области}

\taskhead{Исходные данные и аналитическое решение}
Для варианта 3 заданы $a=2$, $b=3$. Область
$D=\{x\geq0,\ y\geq0,\ x/a+y/b\leq1\}$ --- треугольник с вершинами $(0,0)$, $(a,0)$, $(0,b)$.
\[
I_2=\iint_D(x+y)\,dx\,dy.
\]
\todo{D: выписать пределы интегрирования для своего варианта и вычислить $I_2$ аналитически (пример: main.tex, строки 348–357).}

\taskhead{Замена переменных и якобиан}
% Переход к единичному квадрату: x=a*xi, y=b*(1-xi)*eta, 0<=xi,eta<=1
% (теория: main.tex, строки 225–247; пример: 359–375).
\todo{D: выписать замену для своего варианта, матрицу частных производных, вычислить якобиан $J$; записать преобразованный интеграл по единичному квадрату от $g(\xi,\eta)=f\bigl(x(\xi,\eta),y(\xi,\eta)\bigr)\,|J|$.}

\taskhead{Программный результат и сравнение}
\todo{D: листинг вызовов функций; оба метода применяются к $g$ на единичном квадрате, $h_\xi=h_\eta=1/n$.}
Аналитическое значение $I_2$: \todo{D: подставить}.
\begin{table}[htbp]
\centering\small
\caption{Сравнение численного интеграла с аналитическим значением $I_2$}
\label{tab:1-2}
\begin{tabular}{|l|r|r|r|r|r|}\hline
Метод & $n$ & $h$ & $I_n$ & $\Delta_n$ & $\delta_n$\\\hline
Ср. прямоугольники & 4 & — & — & — & —\\
Ср. прямоугольники & 8 & — & — & — & —\\
Ср. прямоугольники & 16 & — & — & — & —\\
Ср. прямоугольники & 32 & — & — & — & —\\
Ср. прямоугольники & 64 & — & — & — & —\\
Трапеции & 4 & — & — & — & —\\
Трапеции & 8 & — & — & — & —\\
Трапеции & 16 & — & — & — & —\\
Трапеции & 32 & — & — & — & —\\
Трапеции & 64 & — & — & — & —\\
\hline\end{tabular}
\end{table}

\taskhead{Графическая проверка}
\taskfigure[0.6]{figures/1-2-region.pdf}{Область интегрирования — треугольник $D$ (равный масштаб осей)}{fig:1-2-region}
\taskfigure[0.75]{figures/1-2-error.pdf}{Абсолютная ошибка обоих методов от шага сетки (логарифмические оси)}{fig:1-2-error}
\todo{D: пояснить графики.}

\taskhead{Вывод}
\todo{D: вывод о сходимости обоих методов на треугольной области.}


\section{Тройные интегралы и замена координат}
% ======================================================================
% 2.1 Поворот декартовых координат. Ответственный: D — Филиппов Андрей.
%
% Что сдать (main.tex: условие 436–456, порядок 457–487, пример 628–668;
% требования к отчёту 1862–1870):
%   1. Изображения исходной области (x,y,z) и новой (u,v,w), оси и границы.
%   2. Матрица частных производных обратной замены, её определитель, проверка |J|=1.
%   3. Подстановка x,y,z в функцию, упрощение, повторный интеграл по u,v,w
%      с явными пределами и множителем якобиана.
%   4. Аналитическое вычисление последовательным интегрированием.
%   5. Метод средних прямоугольников при n=8,16,32,64; шаги h1=a/n, h2=b/n, h3=c/n.
%   6. Проверка объёма: интеграл от 1 (с якобианом) сравнить с V=abc.
%   7. Таблица значений и ошибок, график ошибки от h=1/n, вывод о сходимости.
%
% Параметры варианта 3 (приложение А): a=2, b=3, c=1, alpha=pi/3.
% Код: solution.py (блок D), запуск python experiments.py 2.1.
% Рисунки: figures/2-1-*.pdf.
% ======================================================================

\subsection{Поворот декартовых координат}

\taskhead{Исходные данные и аналитическое решение}
Для варианта 3 заданы $a=2$, $b=3$, $c=1$, $\alpha=\pi/3$. Рассматривается интеграл
\[
I_{\mathrm{п}}=\iiint_{\Omega_{\mathrm{п}}}(x^2+y^2+z^2)\,dx\,dy\,dz,
\]
\[
\Omega_{\mathrm{п}}=\{0\leq x\cos\alpha+y\sin\alpha\leq a,\quad
0\leq -x\sin\alpha+y\cos\alpha\leq b,\quad 0\leq z\leq c\}.
\]
\todo{D: обосновать выбор координат; выписать поворот и обратную замену с параметрами варианта; составить матрицу частных производных обратной замены и вычислить определитель (проверить $|J|=1$); подставить $x,y,z$ в функцию и упростить; записать повторный интеграл по $u,v,w$; вычислить аналитически (пример: main.tex, строки 628–668).}

\taskhead{Программный результат и сравнение}
\todo{D: листинг вызова своей функции из solution.py; указать шаги $h_1=a/n$, $h_2=b/n$, $h_3=c/n$ для каждой сетки.}
Аналитическое значение $I_{\mathrm{п}}$: \todo{D: подставить}.
На графике используется нормированный шаг $h=1/n$ (main.tex, строки 1989–1990, 2078–2082).
\begin{table}[htbp]
\centering\small
\caption{Сравнение численного интеграла с аналитическим значением $I_{\mathrm{п}}$}
\label{tab:2-1}
\begin{tabular}{|r|r|r|r|r|}\hline
$n$ & $h$ & $I_n$ & $\Delta_n$ & $\delta_n$\\\hline
8 & — & — & — & —\\
16 & — & — & — & —\\
32 & — & — & — & —\\
64 & — & — & — & —\\
\hline\end{tabular}
\end{table}

\taskhead{Проверка объёма}
\todo{D: записать интеграл от 1 по исходной области, преобразовать с учётом якобиана, вычислить численно на всех четырёх сетках и сравнить с $V=abc$.}

\taskhead{Графическая проверка}
\taskfigure[0.55]{figures/2-1-region-old.pdf}{Исходная область в координатах $(x,y,z)$ (равный масштаб осей)}{fig:2-1-region-old}
\taskfigure[0.55]{figures/2-1-region-new.pdf}{Новая область в координатах $(u,v,w)$ (равный масштаб осей)}{fig:2-1-region-new}
\taskfigure[0.75]{figures/2-1-error.pdf}{Абсолютная ошибка от нормированного шага $h=1/n$ (логарифмические оси)}{fig:2-1-error}
\todo{D: пояснить график.}

\taskhead{Вывод}
\todo{D: вывод о сходимости; объяснить поведение ошибки при удвоении $n$.}

% ======================================================================
% 2.2 Цилиндрические координаты. Ответственный: D — Филиппов Андрей.
%
% ЭТА ЖЕ ЗАДАЧА разобрана в report_template.tex (строки 92–200) —
% главный ориентир оформления: блоки, таблица, рисунки, вывод.
%
% Что сдать (main.tex: условие 490–526, порядок 527–553, пример 670–706;
% требования к отчёту 1862–1870):
%   1. Изображение половины цилиндра: радиус, высота, граница y=0.
%   2. Обоснование пределов r, theta, z; замена координат; якобиан
%      вычислением определителя матрицы частных производных.
%   3. Повторный интеграл от g(r,theta,z) с явными пределами; аналитическое решение.
%   4. Метод средних прямоугольников при n=8,16,32,64;
%      шаги h_r=R/n, h_theta=pi/n, h_z=H/n.
%   5. Проверка объёма: интегрировать r (якобиан), сравнить с V=pi*R^2*H/2.
%   6. Таблица значений и ошибок, график ошибки от h=1/n, вывод о сходимости.
%
% Параметры варианта 3 (приложение А): R=1, H=3.
% Код: solution.py (блок D), запуск python experiments.py 2.2.
% Рисунки: figures/2-2-*.pdf.
% ======================================================================

\subsection{Цилиндрические координаты}

\taskhead{Исходные данные и аналитическое решение}
Для варианта 3 заданы $R=1$, $H=3$. Рассматривается интеграл
\[
I_{\mathrm{ц}}=\iiint_{\Omega_{\mathrm{ц}}}(1+x^2+y^2+z)\,dx\,dy\,dz,
\qquad
\Omega_{\mathrm{ц}}=\{x^2+y^2\leq R^2,\ y\geq0,\ 0\leq z\leq H\}.
\]
\todo{D: описать область; обосновать пределы $r,\theta,z$; выписать замену координат; вычислить якобиан через определитель матрицы частных производных; записать функцию $g(r,\theta,z)$ с множителем якобиана и повторный интеграл; вычислить аналитически (пример: main.tex, строки 670–706; оформление: \texttt{report\_template.tex}, строки 95–131).}

\taskhead{Программный результат и сравнение}
\todo{D: листинг вызова своей функции из solution.py (образец: \texttt{report\_template.tex}, строки 138–149); указать шаги $h_r=R/n$, $h_\theta=\pi/n$, $h_z=H/n$.}
Аналитическое значение $I_{\mathrm{ц}}$: \todo{D: подставить}.
\begin{table}[htbp]
\centering\small
\caption{Сравнение численного интеграла с аналитическим значением $I_{\mathrm{ц}}$}
\label{tab:2-2}
\begin{tabular}{|r|r|r|r|r|}\hline
$n$ & $h$ & $I_n$ & $\Delta_n$ & $\delta_n$\\\hline
8 & — & — & — & —\\
16 & — & — & — & —\\
32 & — & — & — & —\\
64 & — & — & — & —\\
\hline\end{tabular}
\end{table}

\taskhead{Проверка объёма}
\todo{D: заменить исходную функцию на 1 (в новых координатах интегрировать $r$, а не 1); вычислить на всех сетках и сравнить с $V=\pi R^2H/2$; пояснить результат (образец: \texttt{report\_template.tex}, строки 173–176).}

\taskhead{Графическая проверка}
\taskfigure[0.6]{figures/2-2-region.pdf}{Половина цилиндра при $R=1$, $H=3$ (равный масштаб осей)}{fig:2-2-region}
\taskfigure[0.75]{figures/2-2-error.pdf}{Абсолютная ошибка от нормированного шага $h=1/n$ (логарифмические оси)}{fig:2-2-error}
\todo{D: пояснить график.}

\taskhead{Вывод}
\todo{D: вывод о сходимости; что подтверждает проверка объёма, а что — нет (образец: \texttt{report\_template.tex}, строки 195–200).}

% ======================================================================
% 2.3 Сферические координаты. Ответственный: D — Филиппов Андрей.
%
% Что сдать (main.tex: условие 555–591, порядок 592–620, пример 708–750;
% требования к отчёту 1862–1870):
%   1. Изображение полушара: радиус, плоскость z=0.
%   2. Обоснование пределов r, theta, phi; матрица частных производных
%      x,y,z по r,theta,phi; определитель; проверка |J|=r^2*sin(phi).
%   3. Повторный интеграл от g(r,theta,phi) с явными пределами; аналитическое решение.
%   4. Метод средних прямоугольников при n=8,16,32,64;
%      шаги h_r=R/n, h_theta=2*pi/n, h_phi=pi/(2n).
%   5. Проверка объёма: интегрировать r^2*sin(phi), сравнить с V=2*pi*R^3/3.
%   6. Таблица значений и ошибок, график ошибки от h=1/n, вывод о сходимости.
%
% Параметры варианта 3 (приложение А): R=1.
% Код: solution.py (блок D), запуск python experiments.py 2.3.
% Рисунки: figures/2-3-*.pdf.
% ======================================================================

\subsection{Сферические координаты}

\taskhead{Исходные данные и аналитическое решение}
Для варианта 3 задан $R=1$. Рассматривается интеграл
\[
I_{\mathrm{с}}=\iiint_{\Omega_{\mathrm{с}}}z^2\,dx\,dy\,dz,
\qquad
\Omega_{\mathrm{с}}=\{x^2+y^2+z^2\leq R^2,\ z\geq0\}.
\]
\todo{D: описать область; обосновать пределы $r,\theta,\varphi$; выписать замену координат; составить матрицу частных производных и вычислить определитель (проверить $|J|=r^2\sin\varphi$); записать функцию $g(r,\theta,\varphi)$ и повторный интеграл; вычислить аналитически (пример: main.tex, строки 708–750).}

\taskhead{Программный результат и сравнение}
\todo{D: листинг вызова своей функции из solution.py; указать шаги $h_r=R/n$, $h_\theta=2\pi/n$, $h_\varphi=\pi/(2n)$.}
Аналитическое значение $I_{\mathrm{с}}$: \todo{D: подставить}.
\begin{table}[htbp]
\centering\small
\caption{Сравнение численного интеграла с аналитическим значением $I_{\mathrm{с}}$}
\label{tab:2-3}
\begin{tabular}{|r|r|r|r|r|}\hline
$n$ & $h$ & $I_n$ & $\Delta_n$ & $\delta_n$\\\hline
8 & — & — & — & —\\
16 & — & — & — & —\\
32 & — & — & — & —\\
64 & — & — & — & —\\
\hline\end{tabular}
\end{table}

\taskhead{Проверка объёма}
\todo{D: заменить исходную функцию на 1 (в новых координатах интегрировать $r^2\sin\varphi$); вычислить на всех сетках и сравнить с $V=2\pi R^3/3$; пояснить результат.}

\taskhead{Графическая проверка}
\taskfigure[0.6]{figures/2-3-region.pdf}{Верхняя половина шара при $R=1$ (равный масштаб осей)}{fig:2-3-region}
\taskfigure[0.75]{figures/2-3-error.pdf}{Абсолютная ошибка от нормированного шага $h=1/n$ (логарифмические оси)}{fig:2-3-error}
\todo{D: пояснить график.}

\taskhead{Вывод}
\todo{D: вывод о сходимости; объяснить поведение ошибки при удвоении $n$.}


\section{Исследование и оптимизация динамической системы}
% ======================================================================
% 3.1 Равновесие и чувствительность. Ответственный: A — Плешнев Арсений.
%
% Что сдать (main.tex: модель и теория 761–812, порядок 935–982,
% разобранный пример 984–1072; требования к отчёту 1871–1879):
%   1. Проверить размерности членов системы модели; объяснить, почему
%      при mu*c <= lambda равновесия нет.
%   2. Найти q* из условия dot q=0; по закону Литтла выразить T(lambda,mu,c);
%      показать, что T зависит от параметров только через запас s.
%   3. Найти T_lambda, T_c, T_mu и дифференциал dT; вычислить в рабочей точке
%      и пояснить смысл; составить уравнение касательной плоскости T=T(lambda,c).
%   4. Производная по направлению (Delta lambda, Delta c); направление нулевого
%      изменения и доказательство постоянства T вдоль него; уравнения линий
%      уровня и их график в плоскости (lambda,c) при фиксированном mu.
%   5. Градиент T; эластичности E_lambda, E_c; вывод: что сильнее влияет
%      на задержку — рост трафика на 1% или рост ресурса на 1%; знаки.
%   6* (по желанию, на оценку не влияет): худший сценарий методом Лагранжа,
%      delta=0.02; сравнить линейный прогноз с точным значением.
%
% Параметры варианта 3 (приложение Б): mu=80, lambda=600, c=8 (рабочая точка);
% a=1, b=1, p=0.03; (Delta lambda, Delta c)=(20, 0.1); mu2=120, N=16.
% ВНИМАНИЕ: a,b,c приложения А (интегралы) и приложения Б (модель) — разные величины.
%
% Аналитику оформить здесь; численная проверка производных — п. 3.5.1.
% Программные значения для таблицы — task_3_1 (python experiments.py 3.1).
% Рисунок линий уровня: figures/3-1-levels.pdf.
% ======================================================================

\subsection{Равновесие и чувствительность}

\taskhead{Модель и исходные данные}
\todo{A: записать систему уравнений модели и обозначения (main.tex, строки 761–812); подставить рабочую точку варианта 3; вычислить запас $s$ и загрузку $\rho$.}

\taskhead{Размерности и условие равновесия}
\todo{A: проверить размерности членов уравнений модели; объяснить, почему при $\mu c\leq\lambda$ система не может прийти в равновесие.}

\taskhead{Равновесная очередь и задержка}
\todo{A: найти $q^*$ из условия $\dot q=0$; по закону Литтла выразить $T(\lambda,\mu,c)$; показать, что $T$ зависит от параметров только через $s$.}

\taskhead{Производные, дифференциал и касательная плоскость}
\todo{A: вывести $T_\lambda,T_c,T_\mu$ и $dT$; вычислить в рабочей точке и пояснить смысл; составить уравнение касательной плоскости к поверхности $T=T(\lambda,c)$ и показать, что оно задаёт линейный прогноз задержки.}

\taskhead{Производная по направлению и линии уровня}
\todo{A: найти производную по направлению $(\Delta\lambda,\Delta c)=(20,\,0.1)$; найти направление нулевого изменения $T$ и доказать, что вдоль него $T$ постоянна точно; записать уравнения линий уровня.}
\taskfigure[0.75]{figures/3-1-levels.pdf}{Линии уровня задержки в плоскости $(\lambda,c)$ при фиксированном $\mu$; отмечена рабочая точка}{fig:3-1-levels}
\todo{A: пояснить график: несколько положительных значений $T$, подписи линий, рабочая точка.}

\taskhead{Градиент и эластичности}
\todo{A: найти $\operatorname{grad}T$; вычислить $E_\lambda,E_c$; сделать вывод, что сильнее влияет на задержку при одинаковом малом относительном изменении, пояснить знаки эластичностей.}

\taskhead{Сопоставление с программным расчётом}
% Программные значения — вызовы своих функций из solution.py (task_3_1);
% численные производные — из п. 3.5.1. Формат таблицы: main.tex, строки 1905–1915.
\todo{A: листинг вызовов функций и запуск python experiments.py 3.1.}
\begin{table}[htbp]
\centering\footnotesize
\setlength{\tabcolsep}{4pt}
\caption{Сравнение аналитических и программных значений (п. 3.1)}
\label{tab:3-1}
\begin{tabular}{|l|r|r|r|r|}\hline
Величина & \shortstack{Аналитическое\\значение} & \shortstack{Программное\\значение} & \shortstack{Абсолютная\\ошибка} & \shortstack{Относительная\\ошибка}\\\hline
$q^*$ & — & — & — & —\\
$T$ & — & — & — & —\\
$T_\lambda$ & — & — & — & —\\
$T_c$ & — & — & — & —\\
$T_\mu$ & — & — & — & —\\
$E_\lambda$ & — & — & — & —\\
$E_c$ & — & — & — & —\\
\hline\end{tabular}
\end{table}

\taskhead{Вывод}
\todo{A: выводы по чувствительности задержки; расхождения и их причины.}

% Пункт 3.1.6* (худший сценарий, метод Лагранжа, delta=0.02) — по желанию,
% на оценку не влияет. Если выполняется: функция Лагранжа, решение системы,
% обоснование максимума, линейный прогноз и сравнение с точным значением.

% ======================================================================
% 3.2 Высшие производные и пределы прогноза. Ответственный: B — Попадюк Евгений.
%
% Что сдать (main.tex: теория 1077–1120, порядок 1121–1170,
% разобранный пример 1172–1276; требования к отчёту 1871–1879):
%   1. Вторые частные производные T(lambda,c); равенство смешанных;
%      определитель матрицы Гессе и объяснение, почему он равен нулю.
%   2. Доказать по индукции формулу n-й производной по lambda;
%      записать ряд Тейлора T(lambda+h), найти его сумму и радиус сходимости.
%   3. При росте трафика h=s/2 вычислить прогнозы Тейлора 1, 2, 3 порядков
%      и их ошибки; сделать вывод.
%   4. По формуле Тейлора 2-го порядка для двух переменных спрогнозировать
%      задержку при (Delta lambda, Delta c) своего варианта; сравнить с точной.
%   5* (по желанию, на оценку не влияет): функция выгоды G, стационарная точка,
%      критерий Сильвестра; объяснить результат.
%
% Параметры варианта 3 (приложение Б): mu=80, lambda=600, c=8;
% (Delta lambda, Delta c)=(20, 0.1); h=s/2; для п. 5*: a=1, b=1, p=0.03.
% Рабочая точка и s — из п. 3.1 (файл 31-equilibrium.tex).
%
% Программные значения для таблиц — task_3_2 (python experiments.py 3.2);
% графики ошибок прогнозов от h — п. 3.5.2 (твой же файл 352-taylor-errors.tex).
% ======================================================================

\subsection{Высшие производные и пределы прогноза}

\taskhead{Исходные данные}
\todo{B: рабочая точка варианта 3 и запас $s$ (из п. 3.1); приращения $(\Delta\lambda,\Delta c)=(20,\,0.1)$; величина $h=s/2$.}

\taskhead{Матрица Гессе}
\todo{B: найти вторые частные производные $T(\lambda,c)$, проверить равенство смешанных; составить матрицу Гессе, вычислить определитель и объяснить, почему он равен нулю (связь с п. 3.1, направление нулевого изменения).}

\taskhead{Ряд Тейлора по трафику}
\todo{B: доказать по индукции формулу $\partial^nT/\partial\lambda^n$; записать ряд Тейлора $T(\lambda+h)$ по степеням $h$, найти его сумму и радиус сходимости (пример: main.tex, строки 1182–1192).}

\taskhead{Прогнозы при росте трафика на $h=s/2$}
\todo{B: вычислить прогнозы $P_1,P_2,P_3$, точное значение задержки и ошибки прогнозов; сделать вывод (пример: main.tex, строки 1193–1208).}

\taskhead{Прогноз по двум переменным}
\todo{B: по формуле Тейлора 2-го порядка спрогнозировать задержку при $(\Delta\lambda,\Delta c)=(20,\,0.1)$; сравнить с точным значением (пример: main.tex, строки 1209–1240).}

\taskhead{Сопоставление с программным расчётом}
% Программные значения — вызовы своих функций из solution.py (task_3_2).
\todo{B: листинг вызовов функций и запуск python experiments.py 3.2.}
\begin{table}[htbp]
\centering\footnotesize
\setlength{\tabcolsep}{4pt}
\caption{Сравнение аналитических и программных значений (п. 3.2)}
\label{tab:3-2}
\begin{tabular}{|l|r|r|r|r|}\hline
Величина & \shortstack{Аналитическое\\значение} & \shortstack{Программное\\значение} & \shortstack{Абсолютная\\ошибка} & \shortstack{Относительная\\ошибка}\\\hline
$P_1$ при $h=s/2$ & — & — & — & —\\
$P_2$ при $h=s/2$ & — & — & — & —\\
$P_3$ при $h=s/2$ & — & — & — & —\\
$P_2$ по двум переменным & — & — & — & —\\
$T$ точное & — & — & — & —\\
\hline\end{tabular}
\end{table}

\taskhead{Вывод}
\todo{B: как порядок прогноза влияет на ошибку; где прогноз перестаёт работать.}

% Пункт 3.2.5* (функция выгоды G, стационарная точка, критерий Сильвестра) —
% по желанию, на оценку не влияет. Если выполняется: производные G, стационарная
% точка, матрица Гессе в ней, обоснованный вывод о характере точки.

% ======================================================================
% 3.3 Динамика регулятора. Ответственный: B — Попадюк Евгений.
%
% Что сдать (main.tex: теория 1281–1315, порядок 1316–1335,
% разобранный пример 1337–1375; требования к отчёту 1871–1879):
%   1. Найти положение равновесия (q*, c*) системы модели.
%   2. Найти матрицу Якоби J правых частей в равновесии, обозначив
%      A=lambda*(1-u0)^2/u0; объяснить, почему d(dot c)/dc = 0.
%   3. Найти собственные числа nu_{1,2}; показать устойчивость при любом k>0;
%      найти порог k*, при превышении которого возникают колебания (флаппинг).
%   4. Для своего lambda вычислить A и k* при u0=0.6; 0.75; 0.9;
%      сделать вывод о поведении системы при u0 -> 1.
%
% Параметры варианта 3 (приложение Б): mu=80, lambda=600.
%
% Программные значения — task_3_3 (python experiments.py 3.3).
% Численный порог флаппинга — п. 3.5.4 (файл 353-simulator.tex, зона C):
% он сравнивается с k*, но расхождение НЕ называется ошибкой аппроксимации —
% результат зависит от критерия обнаружения (main.tex, строки 1652–1661).
% ======================================================================

\subsection{Динамика регулятора}

\taskhead{Исходные данные и положение равновесия}
\todo{B: найти $(q^*,c^*)$ системы модели; вычислить для заданных целевых загрузок (пример: main.tex, строки 1338–1346).}

\taskhead{Матрица Якоби в равновесии}
\todo{B: выписать частные производные правых частей; составить $J$ в равновесии с обозначением $A=\lambda(1-u_0)^2/u_0$; объяснить, почему $\partial\dot c/\partial c=0$ (пример: main.tex, строки 1347–1362).}

\taskhead{Собственные числа и устойчивость}
\todo{B: найти $\nu_{1,2}$; обосновать локальную асимптотическую устойчивость при любом $k>0$; вывести порог $k_*$ появления комплексных собственных чисел (колебаний); пояснить, что этот порог — не порог потери устойчивости.}

\taskhead{Расчёт для трёх целевых загрузок}
\todo{B: вычислить $A$ и $k_*$ при $u_0=0.6;\,0.75;\,0.9$ для своего $\lambda$; сделать вывод о поведении порога при $u_0\to1$.}
\begin{table}[htbp]
\centering\small
\caption{Равновесие и порог колебаний для трёх целевых загрузок}
\label{tab:3-3}
\begin{tabular}{|r|r|r|r|r|}\hline
$u_0$ & $q^*$ & $c^*$ & $A$ & $k_*$\\\hline
0.6 & — & — & — & —\\
0.75 & — & — & — & —\\
0.9 & — & — & — & —\\
\hline\end{tabular}
\end{table}

\taskhead{Сопоставление с программным расчётом}
% Программные значения — вызовы своих функций из solution.py (task_3_3).
\todo{B: листинг вызовов функций и запуск python experiments.py 3.3.}
\begin{table}[htbp]
\centering\footnotesize
\setlength{\tabcolsep}{4pt}
\caption{Сравнение аналитических и программных значений (п. 3.3)}
\label{tab:3-3-check}
\begin{tabular}{|l|r|r|r|r|}\hline
Величина & \shortstack{Аналитическое\\значение} & \shortstack{Программное\\значение} & \shortstack{Абсолютная\\ошибка} & \shortstack{Относительная\\ошибка}\\\hline
$k_*$ при $u_0=0.6$ & — & — & — & —\\
$k_*$ при $u_0=0.75$ & — & — & — & —\\
$k_*$ при $u_0=0.9$ & — & — & — & —\\
\hline\end{tabular}
\end{table}

\taskhead{Вывод}
\todo{B: устойчивость при любом $k>0$; смысл порога $k_*$; поведение при $u_0\to1$.}

% ======================================================================
% 3.4 Оптимизация затрат. Ответственный: A — Плешнев Арсений.
%
% Что сдать (main.tex: теория 1380–1414, порядок 1415–1436,
% разобранный пример 1438–1493; требования к отчёту 1871–1879):
%   1. Объяснить, почему в минимуме производная по любому направлению равна
%      нулю (grad Phi = 0); найти стационарную точку, выразить c_i* через
%      базовую нагрузку R_i=lambda/mu_i (правило квадратного корня).
%   2. Доказать, что точка — минимум (Гессе, выпуклость -> глобальный минимум);
%      вычислить c1* и сравнить затраты первого звена при исходном c и при c1*.
%   3. Найти оптимальную целевую загрузку u0*=lambda/(mu1*c1*) и соответствующий
%      порог k* из п. 3.3.
%   4. При ограничении c1+c2=N: условие оптимальности через производную
%      по направлению (1,-1) и то же условие методом Лагранжа; найти c1,c2
%      численно; проверить сумму и положительность запасов; смысл множителя.
%
% Параметры варианта 3 (приложение Б): lambda=600, mu1=mu=80, mu2=120,
% a=1, b=1, N=16; исходное c=8 (рабочая точка из п. 3.1).
%
% Программные значения — task_3_4 (python experiments.py 3.4);
% градиентный спуск для этой же функции — п. 3.5.5 (файл 353-simulator.tex, зона C).
% ======================================================================

\subsection{Оптимизация затрат}

\taskhead{Исходные данные и функция затрат}
\todo{A: записать $\Phi(c_1,c_2)$ с параметрами варианта 3; допустимая область $s_i=\mu_ic_i-\lambda>0$; базовые нагрузки $R_i=\lambda/\mu_i$.}

\taskhead{Стационарная точка и доказательство минимума}
\todo{A: объяснить, почему $\operatorname{grad}\Phi=0$ в точке минимума; найти стационарную точку и выразить $c_i^*$ через $R_i$; доказать минимум (вторые производные, выпуклость на всей допустимой области -> глобальный характер); вычислить значения (пример: main.tex, строки 1440–1457).}

\taskhead{Затраты и оптимальные настройки регулятора}
\todo{A: сравнить затраты первого звена при исходном $c$ и при $c_1^*$; найти $u_0^*=\lambda/(\mu_1c_1^*)$ и соответствующий порог $k_*$ по формуле п. 3.3.}

\taskhead{Ограниченный общий ресурс $c_1+c_2=N$}
\todo{A: получить условие оптимального распределения из равенства нулю производной по направлению $(1,-1)$; получить то же условие методом Лагранжа; найти $c_1,c_2$ численно (python experiments.py 3.4); проверить ограничение на сумму и положительность запасов; пояснить смысл множителя Лагранжа (пример: main.tex, строки 1467–1493).}

\taskhead{Сопоставление с программным расчётом}
% Программные значения — вызовы своих функций из solution.py (task_3_4).
\todo{A: листинг вызовов функций и запуск python experiments.py 3.4.}
\begin{table}[htbp]
\centering\footnotesize
\setlength{\tabcolsep}{4pt}
\caption{Сравнение аналитических и программных значений (п. 3.4)}
\label{tab:3-4}
\begin{tabular}{|l|r|r|r|r|}\hline
Величина & \shortstack{Аналитическое\\значение} & \shortstack{Программное\\значение} & \shortstack{Абсолютная\\ошибка} & \shortstack{Относительная\\ошибка}\\\hline
$c_1^*$ & — & — & — & —\\
$c_2^*$ & — & — & — & —\\
$\Phi^*$ & — & — & — & —\\
$u_0^*$ & — & — & — & —\\
$k_*$ при $u_0^*$ & — & — & — & —\\
$c_1$ при $c_1+c_2=N$ & — & — & — & —\\
$c_2$ при $c_1+c_2=N$ & — & — & — & —\\
\hline\end{tabular}
\end{table}
% Для распределения при ограничении N аналитического ответа в замкнутой форме
% может не быть: тогда указать способ проверки (main.tex, строки 1911–1915)
% и не называть разность двух приближений ошибкой относительно точного решения.

\taskhead{Вывод}
\todo{A: правило квадратного корня и его смысл; выигрыш в затратах; эффект ограничения ресурса.}

\subsection{Численная проверка}
% ======================================================================
% Эксперимент 3.5.1. Численные производные и оптимальный шаг.
% Ответственный: A — Плешнев Арсений.
% Подраздел 3.5 «Численная проверка» начинается в report.tex; этот файл —
% первый эксперимент подраздела (заголовок subsection не нужен).
%
% Что сдать (main.tex: теория 1504–1516, пункт 1541–1544, ожидаемый результат
% 1563–1565, пример 1583–1624; требования к отчёту 1875–1879):
%   * функции численных частных производных и градиента (односторонние
%     и центральные разности) — в solution.py;
%   * график ошибки оценки dT/dlambda от шага h в логарифмическом масштабе;
%   * найденные по эксперименту оптимальные шаги для обеих схем;
%   * объяснение, почему ошибка не убывает неограниченно (остаток Тейлора
%     растёт с h, ошибка округления растёт как eps/h);
%   * указать точку, диапазон шагов, допуски; сравнить с оценками масштаба
%     оптимального шага из примера.
%
% Параметры варианта 3: рабочая точка из п. 3.1 (mu=80, lambda=600, c=8).
% Запуск: python experiments.py 3.5.1. Рисунок: figures/3-5-1-diff-error.pdf.
% ======================================================================

\taskhead{Эксперимент 3.5.1. Численные производные и оптимальный шаг}

\taskhead{Постановка эксперимента}
\todo{A: описать схемы (односторонняя и центральная разности), точку вычисления, диапазон шагов $h$, точное значение $T_\lambda$ для сравнения; указать допуски.}

\taskhead{Программный результат и сравнение}
\todo{A: листинг вызова своих функций из solution.py и запуск python experiments.py 3.5.1.}
\taskfigure[0.75]{figures/3-5-1-diff-error.pdf}{Ошибка оценки $\partial T/\partial\lambda$ от шага $h$ для обеих схем (логарифмические оси)}{fig:3-5-1}
\begin{table}[htbp]
\centering\small
\caption{Оптимальные шаги численного дифференцирования}
\label{tab:3-5-1}
\begin{tabular}{|l|r|r|}\hline
Схема & Оптимальный шаг $h$ & Минимальная ошибка\\\hline
Односторонняя разность & — & —\\
Центральная разность & — & —\\
\hline\end{tabular}
\end{table}

\taskhead{Вывод}
\todo{A: почему ошибка не убывает неограниченно при уменьшении $h$; сравнить найденные шаги с оценками масштаба $s\sqrt{\varepsilon}$ и $s\varepsilon^{1/3}$ (main.tex, строки 1616–1624).}

% ======================================================================
% Эксперимент 3.5.2. Ошибки прогнозов Тейлора в зависимости от h.
% Ответственный: B — Попадюк Евгений.
%
% Что сдать (main.tex: пункт 1545–1546, ожидаемый результат 1566–1567;
% прогнозы — из п. 3.2, файл 32-taylor.tex):
%   * график ошибок прогнозов P1, P2, P3 из п. 3.2.3 от приращения трафика h;
%   * показать, что прогноз перестаёт работать при h -> s;
%   * указать диапазон h, рабочую точку, точное значение T(lambda+h)=1/(s-h).
%
% Параметры варианта 3: рабочая точка из п. 3.1; s — оттуда же.
% Запуск: python experiments.py 3.5.2. Рисунок: figures/3-5-2-taylor-error.pdf.
% ======================================================================

\taskhead{Эксперимент 3.5.2. Ошибки прогнозов Тейлора}

\taskhead{Постановка эксперимента}
\todo{B: описать, какие прогнозы сравниваются (порядки 1–3 из п. 3.2), диапазон приращений $h$ вплоть до $s$, точное значение для сравнения.}

\taskhead{Программный результат и сравнение}
\todo{B: листинг вызова своих функций из solution.py и запуск python experiments.py 3.5.2.}
\taskfigure[0.75]{figures/3-5-2-taylor-error.pdf}{Ошибки прогнозов Тейлора 1–3 порядков от приращения трафика $h$}{fig:3-5-2}
\todo{B: пояснить график: как ведёт себя ошибка каждого порядка при росте $h$ и что происходит при $h\to s$.}

\taskhead{Вывод}
\todo{B: пределы применимости прогноза фиксированного порядка; связь с радиусом сходимости из п. 3.2.}

% ======================================================================
% Эксперименты 3.5.3–3.5.5. Ответственный: C — Дучин Кирилл.
%
% Что сдать (main.tex: теория 1518–1536, пункты 1547–1557, ожидаемый результат
% 1568–1575, пример 1626–1699; требования к отчёту 1875–1879):
%   3.5.3  Симулятор системы методом Эйлера; найти шаг Delta t потери
%          устойчивости; сравнить с условием |1+Delta t*nu|<1 для собственных
%          чисел из п. 3.3; траектории q(t), c(t) при указанных начальных условиях.
%   3.5.4  Численно найти порог флаппинга для u0=0.6; 0.75; 0.9 и сравнить с k*;
%          описать критерий обнаружения; проверить устойчивость наблюдения
%          при уменьшении временного шага.
%   3.5.5  Градиентный спуск для Phi(c1,c2) из п. 3.4 со стартовой точки (10,10)
%          при постоянном и убывающем шаге; число итераций, точка, значение
%          функции; траектория на фоне линий уровня; сравнение с аналитикой.
%          ОБРАЗЕЦ ОФОРМЛЕНИЯ ЭТОГО ЭКСПЕРИМЕНТА — report_template.tex, 203–336.
%
% Параметры варианта 3 (приложение Б): lambda=600, mu=80, mu2=120, a=1, b=1.
% NB (main.tex, строки 1692–1699): перед запуском проверить допустимость (10,10):
% mu_i*c_i > lambda; если недопустима — использовать (R1+1, R2+1) и указать замену.
% Зафиксировать начальные условия симуляции, длительность и критерий обнаружения
% колебаний. Численный градиент — переиспользовать функции разностей из п. 3.5.1
% (блок A в solution.py), не дублировать.
%
% Запуск: python experiments.py 3.5.3 / 3.5.4 / 3.5.5.
% Рисунки: figures/3-5-3-trajectories.pdf, figures/3-5-5-paths.pdf.
% ======================================================================

\taskhead{Эксперимент 3.5.3. Моделирование методом Эйлера и граница устойчивости}

\taskhead{Постановка эксперимента}
\todo{C: описать симулятор (уравнения, начальные условия, длительность, перебираемые $\Delta t$); указать, как фиксировалась потеря устойчивости.}

\taskhead{Программный результат и сравнение}
\todo{C: листинг вызова своих функций из solution.py и запуск python experiments.py 3.5.3.}
\taskfigure[0.85]{figures/3-5-3-trajectories.pdf}{Траектории $q(t)$ и $c(t)$ при выбранных начальных условиях и настройках}{fig:3-5-3}
\begin{table}[htbp]
\centering\small
\caption{Граница устойчивости по временному шагу: численный порог и линейный критерий}
\label{tab:3-5-3}
\begin{tabular}{|r|r|r|r|r|}\hline
$u_0$ & $k$ & $\Delta t$ (числ. порог) & $2/|\nu|$ (критерий) & Расхождение\\\hline
— & — & — & — & —\\
\hline\end{tabular}
\end{table}
% Расхождение численного порога и линейного критерия — не «ошибка»:
% критерий относится к линейному приближению вблизи равновесия
% (main.tex, строки 1523–1527). Пояснить причину расхождения.
\todo{C: пояснить график и таблицу; сравнить с условием $|1+\Delta t\,\nu|<1$ для $\nu$ из п. 3.3.}

\taskhead{Эксперимент 3.5.4. Численный порог флаппинга}

\taskhead{Постановка эксперимента}
\todo{C: описать критерий обнаружения колебаний (начальные условия, длительность, допуск, перебор $k$); одно пересечение уровня $c^*$ не доказывает колебаний — анализировать повторные смены направления отклонения (main.tex, строки 1652–1661).}

\taskhead{Программный результат и сравнение}
\todo{C: листинг вызовов и запуск python experiments.py 3.5.4.}
\begin{table}[htbp]
\centering\small
\caption{Пороги колебательного режима: аналитический $k_*$ и численный порог}
\label{tab:3-5-4}
\begin{tabular}{|r|r|r|r|}\hline
$u_0$ & $k_*$ (п. 3.3) & Численный порог & Расхождение\\\hline
0.6 & — & — & —\\
0.75 & — & — & —\\
0.9 & — & — & —\\
\hline\end{tabular}
\end{table}
\todo{C: пояснить расхождения (зависимость от начальных условий, шага, длительности, допуска); проверить устойчивость наблюдения при уменьшении $\Delta t$.}

\taskhead{Эксперимент 3.5.5. Градиентный спуск для функции затрат}

\taskhead{Постановка эксперимента}
\todo{C: проверить допустимость $(10,10)$ для варианта 3 (при нарушении — $(R_1+1,R_2+1)$ и указать замену); описать численный градиент (схема, шаг $h_d$), два правила шага $\eta_j$ (постоянный и убывающий — выбрать и записать), критерий остановки, предел итераций, проверку допустимости на каждом шаге (образец: \texttt{report\_template.tex}, строки 237–261).}

\taskhead{Программный результат и сравнение}
\todo{C: листинг вызовов и запуск python experiments.py 3.5.5.}
\begin{table}[htbp]
\centering\small
\caption{Результаты градиентного спуска}
\label{tab:3-5-5}
\begin{tabular}{|l|r|r|r|r|r|}\hline
Шаг & Итерации & $\widehat c_1$ & $\widehat c_2$ & $\Phi(\widehat c)$ & $\Delta_c$\\\hline
Постоянный & — & — & — & — & —\\
Убывающий & — & — & — & — & —\\
\hline\end{tabular}
\end{table}
\taskfigure[0.85]{figures/3-5-5-paths.pdf}{Траектории спуска на фоне линий уровня затрат; отмечены начальная точка и аналитический минимум}{fig:3-5-5}
\todo{C: сравнить с аналитическим минимумом из п. 3.4 ($\Delta_c$, $\delta_c$ — определения в \texttt{report\_template.tex}, строки 274–279); ошибки затрат; допустимость всех точек.}

\taskhead{Вывод}
\todo{C: сравнение правил шага (осторожно: вывод относится только к выбранным правилам и начальной точке); причины расхождений.}

% ======================================================================
% 3.6 Настройка в реалистичных условиях. Ответственный: C — Дучин Кирилл.
%
% Что сдать (main.tex: теория 1702–1734, порядок 1736–1760, пример 1762–1831;
% требования к отчёту 1880–1884):
%   1. Явно задать: функцию трафика lambda(t) (суточная синусоида + всплеск),
%      значение задержки tau, уравнение регулятора с задержкой, начальную
%      историю на [-tau, 0], T_sim, Delta t и способ учёта задержки
%      (целое m=tau/Delta t или интерполяция).
%   2. Вычислить средние затраты за сутки модельного времени (формула
%      трапеций по L_n = a*c_n + b*q_n).
%   3. Найти минимум средних затрат градиентным спуском с численным градиентом,
%      начав с u0* и k* из п. 3.3–3.4; линии уровня на плоскости (u0,k)
%      и траектория спуска.
%   4. Сравнить найденные настройки с аналитическими и ответить на сквозной
%      вопрос работы; объяснить ограничения выводов.
%
% Параметры варианта 3 (приложение Б): lambda=600, mu=80, a=1, b=1.
% ВАЖНО (main.tex, строки 1825–1831): u0* — лишь исходная настройка, k* — порог
% колебаний модели без задержки, а не минимум затрат; выпуклость целевой функции
% не установлена; шаг Delta t, пригодный без задержки, не переносится автоматически
% на систему с задержкой — устойчивость и точность проверить отдельно.
% Во всех прогонах для численного градиента — одинаковые трафик, история, сетка.
%
% Запуск: python experiments.py 3.6.
% Рисунки: figures/3-6-levels.pdf, figures/3-6-dynamics.pdf.
% ======================================================================

\subsection{Настройка в реалистичных условиях}

\taskhead{Постановка эксперимента}
\todo{C: явно выписать $\lambda(t)$ (амплитуда суточного колебания, всплеск: интервал и величина), значение $\tau$, уравнение регулятора с задержкой, начальную историю, $T_{\mathrm{sim}}$, $\Delta t$ и проверку его устойчивости/точности; способ учёта задержки.}

\taskhead{Средние затраты и численный градиент}
\todo{C: описать вычисление $\overline\Phi(u_0,k)$ (формула трапеций по $L_n=ac_n+bq_n$); численный градиент отдельными прогонами с одинаковыми трафиком/историей/сеткой; шаги $\varepsilon_u,\varepsilon_k$; критерий остановки; допустимость $0<u_0<1$, $k>0$; что делать с неустойчивой траекторией.}

\taskhead{Результаты оптимизации}
\todo{C: листинг вызовов и запуск python experiments.py 3.6; старт с $u_0^*,k_*$ из п. 3.3–3.4.}
\begin{table}[htbp]
\centering\small
\caption{Исходные и найденные настройки регулятора}
\label{tab:3-6}
\begin{tabular}{|l|r|r|r|}\hline
Настройки & $u_0$ & $k$ & $\overline\Phi$\\\hline
Исходные (из п. 3.3–3.4) & — & — & —\\
Найденные спуском & — & — & —\\
\hline\end{tabular}
\end{table}
\taskfigure[0.8]{figures/3-6-levels.pdf}{Линии уровня средних затрат в плоскости $(u_0,k)$ и траектория спуска}{fig:3-6-levels}
\taskfigure[0.85]{figures/3-6-dynamics.pdf}{Трафик $\lambda(t)$, очередь $q(t)$ и ресурс $c(t)$ за сутки модельного времени}{fig:3-6-dynamics}
\todo{C: пояснить графики; указать параметры численного градиента и критерий остановки.}

\taskhead{Вывод}
\todo{C: сравнить найденные настройки с аналитическими; ответить на сквозной вопрос работы (как настраивать регулятор, чтобы минимизировать затраты при сохранении устойчивости); ограничения выводов: локальный характер минимума, зависимость от выбранного трафика и $\tau$.}


\section*{Проверка реализации}
% ======================================================================
% Проверка реализации (таблица). Ведёт D — Филиппов Андрей.
% Каждый владелец присылает строки за свой блок (PR меняет только свою строку).
%
% Требования (main.tex, строки 1885–1891): столбцы «Задача», «Команда запуска»,
% «Проверяемая величина», «Результат проверки»; включить сравнение интегралов
% и объёмов с аналитическими значениями и проверки численных методов из п. 3.5.
% В результате — конкретное расхождение или выполненное условие.
% Полный журнал вывода программы вместо таблицы не приводится.
% Образец: report_template.tex, строки 313–336.
% ======================================================================

\begin{table}[htbp]
\centering\small
\caption{Проверка реализации}
\label{tab:checks}
\begin{tabular}{|p{1.4cm}|p{3.9cm}|p{4.6cm}|p{4.6cm}|}\hline
Задача & Команда запуска & Проверяемая величина & Результат проверки\\\hline
1.1 & \texttt{python experiments.py 1.1} & Интеграл $I_1$: расхождение с аналитическим значением на сетках $4,\dots,64$ (оба метода) & —\\\hline
1.2 & \texttt{python experiments.py 1.2} & Интеграл $I_2$: расхождение с аналитическим значением на сетках $4,\dots,64$ (оба метода) & —\\\hline
2.1 & \texttt{python experiments.py 2.1} & Интеграл $I_{\mathrm{п}}$ и объём: расхождение с аналитикой и $V=abc$ & —\\\hline
2.2 & \texttt{python experiments.py 2.2} & Интеграл $I_{\mathrm{ц}}$ и объём: расхождение с аналитикой и $V=\pi R^2H/2$ & —\\\hline
2.3 & \texttt{python experiments.py 2.3} & Интеграл $I_{\mathrm{с}}$ и объём: расхождение с аналитикой и $V=2\pi R^3/3$ & —\\\hline
3.1–3.4 & \texttt{python experiments.py 3.1} и т.д. & Значения формул: расхождение аналитического и программного расчёта & —\\\hline
3.5.1 & \texttt{python experiments.py 3.5.1} & Оптимальные шаги разностей; минимум ошибки на найденном шаге & —\\\hline
3.5.2 & \texttt{python experiments.py 3.5.2} & Ошибки прогнозов $P_1$–$P_3$ (п. 3.2.3) от $h$: рост при $h\to s$ & —\\\hline
3.5.3 & \texttt{python experiments.py 3.5.3} & Граница устойчивости по $\Delta t$: сравнение с линейным критерием & —\\\hline
3.5.4 & \texttt{python experiments.py 3.5.4} & Пороги флаппинга: сравнение с $k_*$; устойчивость наблюдения при уменьшении шага & —\\\hline
3.5.5 & \texttt{python experiments.py 3.5.5} & Координаты минимума ($\Delta_c$); допустимость точек; невозрастание затрат & —\\\hline
3.6 & \texttt{python experiments.py 3.6} & Средние затраты: сходимость при уменьшении $\Delta t$ & —\\\hline
\end{tabular}
\end{table}


\section*{Заключение по лабораторной работе}
% ======================================================================
% Заключение. Ответственный: A — Плешнев Арсений (собирает выводы от всех).
% Требования (main.tex, строки 1892–1897): основные результаты, расхождения
% и их причины; влияние шага сетки, временного шага и шага оптимизации;
% какие выводы подтверждены аналитически, а какие получены только для
% выбранного численного эксперимента; список внешних источников.
% Образец: report_template.tex, строки 338–351.
% ======================================================================

\todo{A: основные результаты по разделам 1–3 (каждый владелец присылает 2–3 предложения по своей части).}

\todo{A: обнаруженные расхождения и их причины (погрешность квадратур, критерии остановки, машинное округление, критерии обнаружения колебаний).}

\todo{A: влияние шага сетки, временного шага и шага оптимизации на результат.}

\todo{A: отдельно указать, какие выводы подтверждены аналитически, а какие получены только для выбранного численного эксперимента.}

\taskhead{Источник исходных данных}
Условие лабораторной работы \texttt{main.tex} (репозиторий преподавателя),
вариант 3: приложение А — параметры интегралов, приложение Б — параметры модели.

\todo{A: список использованных внешних источников, если были.}

\end{document}
